% ============================================================
%  Week 5 — Randomization and Back-Door Adjustment
%  Compile standalone:  pdflatex chapter5_new.tex
%  Compile via master:  pdflatex causal_inference_master.tex
% ============================================================
\documentclass[master.tex]{subfiles}

\begin{document}

% ------------------------------------------------------------
%  Title block (standalone) / chapter heading (master)
% ------------------------------------------------------------
\ifSubfilesClassLoaded{%
  \begin{center}
    {\LARGE\bfseries\color{isublue}
     Chapter 5: Randomization and Back-Door Adjustment}\\[6pt]
    {\large Theory and Methods in Causal Inference}\\[4pt]
    {\normalsize Department of Statistics, Iowa State University}
  \end{center}
  \bigskip
  \hrule height 1.2pt \relax
  \smallskip
}{%
  \chapter{Randomization and Back-Door Adjustment}
}

% ------------------------------------------------------------
%  Learning objectives
% ------------------------------------------------------------
\begin{tcolorbox}[defstyle, title={Learning Objectives}]
By the end of this lecture, students should be able to:
\begin{enumerate}
  \item Explain in do-calculus terms why a randomized experiment makes
    $f(y \mid \doop(T{=}t)) = f(y \mid T{=}t)$, and identify the
    graphical feature that produces this equality.
  \item State the ignorability (unconfoundedness) assumption in both
    potential-outcomes and graphical language, and distinguish weak from
    strong ignorability.
  \item Apply the regression-adjusted estimator to reduce variance in a
    randomized experiment, and explain why consistency holds even under
    model misspecification.
  \item Derive and compute the standardization (g-formula) estimator
    of the ATE from a contingency table or regression output.
  \item Explain the method of stratification (subclassification) and
    describe when and why it removes confounding bias.
  \item Identify Simpson's paradox in a numerical example and resolve
    it using the back-door criterion.
  \item Articulate why propensity score methods are needed when $X$ is
    high-dimensional, motivating Chapter~6.
\end{enumerate}
\end{tcolorbox}


% ============================================================
\section{Randomized Experiments}
\label{w5:sec:rct}
% ============================================================

\subsection{The Do-Calculus of Randomization}

The defining feature of a randomized experiment is that the analyst
\emph{sets} the treatment for each unit, independently of all background
variables.  In do-calculus terms, the observed data come from the
mutilated distribution $f(y \mid \doop(T{=}t))$ rather than the
observational distribution $f(y \mid T{=}t)$.

\begin{tcolorbox}[defstyle, title={The Fundamental Equality of Randomized Experiments}]
In a completely randomized experiment, $T$ is assigned independently of
all pre-treatment variables.  In the DAG, this means there are
\emph{no arrows into $T$}: the treatment node has no parents.
By the do-calculus (Rule~2 applied to the trivial graph
$G_{\overline{T}}$):
\begin{equation}
  f\!\bigl(y \mid \doop(T{=}t)\bigr) = f(y \mid T{=}t).
  \label{w5:eq:fund-eq}
\end{equation}
In a randomized experiment, observing $T = t$ is the same as intervening
to set $T = t$.
\end{tcolorbox}

\noindent\textbf{Graphical argument.}  In the observational DAG, there
are arrows into $T$ from observed covariates $X$ and unobserved
confounders $U$, creating back-door paths from $T$ to $Y$.
Randomization \emph{physically} severs these arrows: assignment is
determined by a coin flip, not by $X$ or $U$.  The mutilated graph
$G_{\overline{T}}$ is, in a randomized experiment, the \emph{actual}
data-generating graph.  There are no back-door paths to block, because
none exist.

\begin{figure}[h]
\centering
\begin{tikzpicture}[
  node distance = 12mm and 18mm,
  obsnode/.style  = {circle, draw=accent, line width=1pt,
                     minimum size=9mm, font=\small, fill=defbg},
  unobsnode/.style= {circle, draw=backdoorred, line width=1pt, dashed,
                     minimum size=9mm, font=\small, fill=warnbg},
  coinnode/.style = {rectangle, rounded corners=3pt,
                     draw=causalgreen, line width=1pt,
                     minimum width=16mm, minimum height=8mm,
                     font=\scriptsize, fill=green!8,
                     align=center},
  obsarr/.style   = {-{Stealth[length=4pt]}, line width=0.9pt,
                     color=accent},
  redarr/.style   = {-{Stealth[length=4pt]}, line width=0.9pt,
                     color=backdoorred, dashed},
  sevarr/.style   = {-{Stealth[length=4pt]}, line width=0.9pt,
                     color=causalgreen},
  lbl/.style      = {font=\footnotesize\color{isublue}}
]
%% ── Left panel: Observational DAG ───────────────────────────────────
\begin{scope}[xshift=0cm]
  \node[obsnode]                          (T1) {$T$};
  \node[obsnode,  right=18mm of T1]       (Y1) {$Y$};
  \node[unobsnode,above=12mm of T1, xshift=9mm] (U1) {$U$};
  \node[obsnode,  below=12mm of T1, xshift=9mm] (X1) {$X$};

  \draw[redarr]  (U1) -- (T1);
  \draw[redarr]  (U1) -- (Y1);
  \draw[obsarr]  (X1) -- (T1);
  \draw[obsarr]  (X1) -- (Y1);
  \draw[obsarr]  (T1) -- (Y1);

  \node[lbl, below=22mm of T1, xshift=9mm, align=center]
    {\textbf{Observational DAG}\\[2pt]
     Arrows into $T$ from $X$ and $U$\\
     create back-door paths
     $T \leftarrow U \rightarrow Y$};
\end{scope}

%% ── Arrow between panels ────────────────────────────────────────────
\node at (6.0cm, 0.3cm)
  {\Large$\xrightarrow{\text{randomize}}$};

%% ── Right panel: Randomized DAG ─────────────────────────────────────
\begin{scope}[xshift=8.0cm]
  \node[obsnode]                          (T2) {$T$};
  \node[obsnode,  right=18mm of T2]       (Y2) {$Y$};
  \node[unobsnode,above=12mm of T2, xshift=9mm] (U2) {$U$};
  \node[obsnode,  below=12mm of T2, xshift=9mm] (X2) {$X$};
  \node[coinnode, left=14mm of T2]        (C2) {coin\\[-1pt]flip};

  %% Severed arrows drawn as light grey stubs with a cross
  \draw[color=lightgray, line width=0.9pt, dashed]  (U2) -- (T2);
  \draw[color=lightgray, line width=0.9pt, dashed]  (X2) -- (T2);
  %% Mark the severance with a small X at the midpoint
  \node[font=\small\bfseries, color=backdoorred]
    at ($(U2)!0.52!(T2)$) {$\times$};
  \node[font=\small\bfseries, color=backdoorred]
    at ($(X2)!0.52!(T2)$) {$\times$};

  %% Remaining active arrows
  \draw[redarr]  (U2) -- (Y2);
  \draw[obsarr]  (X2) -- (Y2);
  \draw[obsarr]  (T2) -- (Y2);
  \draw[sevarr]  (C2) -- (T2);

  \node[lbl, below=22mm of T2, xshift=9mm, align=center]
    {\textbf{Randomized DAG}\ ($G_{\overline{T}}$)\\[2pt]
     Arrows into $T$ severed ($\times$)\\
     $T$ determined by coin flip only};
\end{scope}
\end{tikzpicture}
\caption{Randomization as graph surgery.  \textbf{Left}: in the
observational DAG, the unobserved confounder $U$ and observed
covariate $X$ both have arrows into $T$, creating back-door paths
$T \leftarrow U \rightarrow Y$ and $T \leftarrow X \rightarrow Y$.
\textbf{Right}: randomization physically severs all arrows into $T$
(marked $\times$) and replaces them with an exogenous coin flip.  The
resulting graph is the mutilated graph $G_{\overline{T}}$, in which no
back-door paths from $T$ to $Y$ exist.}
\label{w5:fig:randomization-dag}
\end{figure}

\medskip\noindent\textbf{Potential outcomes statement.}  In the potential
outcomes language, randomization implies:
\[
  \bigl(Y(0),\, Y(1)\bigr) \;\indep\; T.
\]
The potential outcomes are independent of assignment
\emph{unconditionally} --- no covariate adjustment is needed.  This is
the strongest possible version of ignorability
(Definition~\ref{w5:defn:ignorability}).

\subsection{Estimation of the ATE in a Randomized Experiment}

\begin{lemma}{Identification under Complete Randomization}
\label{w5:prop:rct-id}
Under complete randomization,
\begin{equation}
  \E\bigl[Y(t)\bigr] \;=\; \E\bigl[Y \mid T{=}t\bigr],
  \qquad t \in \{0,1\}.
  \label{w5:eq:rct-id}
\end{equation}
\end{lemma}

\begin{proof}
By Theorem~\ref{w4:prop:po-do}, $Y(t)$ has the same distribution as
$Y$ under the interventional distribution $f(y \mid \doop(T{=}t))$, so
$\E[Y(t)] = \E[Y \mid \doop(T{=}t)]$.  In a completely randomized
experiment, the treatment node has no parents in the DAG, so
equation~\eqref{w5:eq:fund-eq} gives
$f(y \mid \doop(T{=}t)) = f(y \mid T{=}t)$.
Therefore $\E[Y(t)] = \E[Y \mid T{=}t]$.
\end{proof}

\noindent Under complete randomization, the ATE is estimated by the
\emph{difference-in-means} estimator:
\[
  \hat\tau_{\mathrm{DIM}}
  = \bar Y_1 - \bar Y_0
  = \frac{1}{n_1}\sum_{i:\, T_i=1} Y_i
  - \frac{1}{n_0}\sum_{i:\, T_i=0} Y_i,
\]
where $n_1$ and $n_0$ are the numbers of treated and control units.

\begin{thm}{Unbiasedness and Variance of $\hat\tau_{\mathrm{DIM}}$}
\label{w5:prop:dim}
Let $\sigma_t^2 = \mathrm{Var}(Y(t))$ for $t \in \{0,1\}$.
Under complete randomization:
\begin{align}
  \E\bigl[\hat\tau_{\mathrm{DIM}}\bigr]
  &= \E[Y(1)] - \E[Y(0)] = \tau_{\mathrm{ATE}},
  \label{w5:eq:dim-mean} \\[4pt]
  \mathrm{Var}\bigl(\hat\tau_{\mathrm{DIM}}\bigr)
  &= \frac{\sigma_1^2}{n_1} + \frac{\sigma_0^2}{n_0}.
  \label{w5:eq:dim-var}
\end{align}
Equation~\eqref{w5:eq:dim-var} is the \emph{Neyman variance}
\citep{neyman1923application}.
\end{thm}

\begin{proof}
\textit{(Unbiasedness.)}
By the consistency assumption, $Y_i = Y_i(T_i)$, so
$\E[Y_i \mid T_i{=}t] = \E[Y_i(t) \mid T_i{=}t]$.
By Lemma~\ref{w5:prop:rct-id},
$\E[Y_i(t) \mid T_i{=}t] = \E[Y \mid T{=}t] = \E[Y(t)]$.
Hence $\E[\bar Y_t] = \E[Y(t)]$ for each arm, and
\[
  \E\bigl[\hat\tau_{\mathrm{DIM}}\bigr]
  = \E[\bar Y_1] - \E[\bar Y_0]
  = \E[Y(1)] - \E[Y(0)] = \tau_{\mathrm{ATE}}.
\]

\smallskip\noindent\textit{(Variance.)}
Under complete randomization, treatment assignments are independent
across units, so $\bar Y_1$ and $\bar Y_0$ are independent.
Within arm $t$, the $n_t$ observed outcomes are i.i.d.\ draws from the
distribution of $Y(t)$ (by Lemma~\ref{w5:prop:rct-id} and
consistency), giving $\mathrm{Var}(\bar Y_t) = \sigma_t^2 / n_t$.
Independence of the two arms then yields
\[
  \mathrm{Var}\bigl(\hat\tau_{\mathrm{DIM}}\bigr)
  = \mathrm{Var}(\bar Y_1) + \mathrm{Var}(\bar Y_0)
  = \frac{\sigma_1^2}{n_1} + \frac{\sigma_0^2}{n_0}. \qedhere
\]
\end{proof}

\subsection{Fisher's Randomization Inference vs.\ Neyman's Repeated Sampling}

Two distinct inferential frameworks apply to randomized experiments.

\medskip
\renewcommand{\arraystretch}{1.4}
\begin{tabular}{p{2.8cm} p{5.5cm} p{5.5cm}}
\toprule
& \textbf{Fisher} & \textbf{Neyman} \\
\midrule
Target
  & Sharp null: $Y_i(1) = Y_i(0)$ for all $i$
  & $\tau_{\mathrm{ATE}} = \E[Y(1) - Y(0)]$ \\
Randomness
  & From the randomization alone;
    potential outcomes are \emph{fixed}
  & From sampling; potential outcomes are random variables \\
Inference
  & Exact $p$-values via permutation distribution
  & Point estimates and confidence intervals \\
Estimand
  & No estimation; only testing
  & $\tau_{\mathrm{ATE}}$ \\
\bottomrule
\end{tabular}
\renewcommand{\arraystretch}{1}
\medskip

In this course we follow the Neyman framework: potential outcomes are
treated as random variables on a super-population, and the goal is
estimation of $\tau_{\mathrm{ATE}}$ (or $\tau_{\mathrm{ATT}}$) with
valid standard errors.

\begin{remark}[Randomization Licenses Inference Without Population Assumptions]
\label{w5:rem:finite-pop}
A fundamental advantage of randomization that deserves explicit emphasis
is that it licenses causal inference \emph{without any assumption about
a population model}.  In Fisher's design-based framework, the
potential outcomes $\{Y_i(0), Y_i(1)\}$ are treated as fixed numbers
attached to the $n$ units in front of the researcher.  The only source
of randomness is the assignment vector $\mathbf{T} = (T_1,\ldots,T_n)$,
which is drawn uniformly at random from the set of all
$\binom{n}{n_1}$ possible assignments.  This is a finite, enumerable
set, and probability statements refer to this mechanism --- not to
repeated draws from a population.

The practical consequence is striking: the i.i.d.\ assumption, which
is the standard foundation of almost every statistical procedure, is
not required.  Units need not be a random sample from any population;
they can be a convenience sample, a census, or a small targeted group.
The randomization itself --- the physical act of assigning treatment by
coin flip or lottery --- is what justifies inference.  Observational
methods, by contrast, must invoke either an i.i.d.\ sampling assumption
or a model-based superpopulation framework, because there is no
analogous physical mechanism to anchor probability statements.  This is
one of the deepest reasons why a well-conducted randomized experiment
is considered more credible than observational analysis: its inferential
validity rests on design, not on assumptions about the world.
\end{remark}

\begin{remark}[Fisher's Sharp Null: Unidentifiable Effects, Exactly Testable Hypothesis]
\label{w5:rem:sharp-null}
A seemingly paradoxical feature of the Fisher framework is that
individual causal effects $\tau_i = Y_i(1) - Y_i(0)$ are
\emph{fundamentally unidentifiable} for any unit --- only one of the
two potential outcomes is ever observed, a constraint known as the
fundamental problem of causal inference --- yet the sharp null
hypothesis
\[
  H_0:\; Y_i(1) = Y_i(0) \;\text{ for all } i = 1,\ldots,n
\]
is \emph{exactly testable} via randomization.

The resolution of the paradox is that the sharp null completely
specifies the joint potential outcome table.  Under $H_0$, every unit's
two potential outcomes are equal, so the observed outcome $Y_i =
Y_i(T_i)$ equals $Y_i(1) = Y_i(0)$ regardless of the assignment.  The
full vector $(Y_1(0), Y_1(1), \ldots, Y_n(0), Y_n(1))$ is therefore
known from the observed data under $H_0$, even though only $n$ of the
$2n$ entries are observed.  The permutation distribution of
$\hat\tau_{\mathrm{DIM}}$ under all $\binom{n}{n_1}$ possible
assignments is then completely determined, and the $p$-value is the
fraction of assignments giving a statistic as extreme as the observed
one --- no estimation step required.

Individual effects are untestable because any specific value of
$\tau_i$ would be a statement about one unobserved potential outcome
of one unit --- a claim no amount of data can verify.  The sharp null
is testable because it is a \emph{joint} restriction on all $n$ units
that, when true, makes the entire data vector fully observed under any
assignment.  This illustrates a broader principle: untestability of
individual-level quantities does not imply untestability of
population-level or joint hypotheses derived from them.
\end{remark}


\begin{remark}[Further Reading on the Randomization Approach]
\label{w5:rem:randomization-reading}
The finite-population, design-based perspective on causal inference
developed in this section has a rich modern literature.
\citet{ding2023first} provides a comprehensive graduate-level treatment
of the randomization framework, covering CRE, stratified experiments,
regression adjustment, and beyond.
The asymptotic theory underlying large-sample confidence intervals for
$\hat\tau_{\mathrm{DIM}}$ --- specifically, a finite-population central
limit theorem that requires neither i.i.d.\ observations nor a
superpopulation model --- is worked out rigorously by
\citet{li2017general}.
For inference about \emph{heterogeneity} of treatment effects across
units (rather than merely the average), \citet{ding2016randomization}
develop randomization-based tests and confidence intervals under the
Neyman--Rubin model; their approach makes no parametric assumptions and
is directly applicable to completely randomized experiments of the type
studied here.
\end{remark}

% ============================================================
\section{Ignorability}
\label{w5:sec:ignorability}
% ============================================================

\subsection{Two Routes to the Same Estimand}

Section~\ref{w5:sec:rct} established that a randomized experiment achieves the
fundamental equality $f(y \mid \doop(T{=}t)) = f(y \mid T{=}t)$ and
that the difference-in-means estimator is unbiased for the ATE.  The
mechanism is graphical: randomization physically severs every arrow into
$T$, so no back-door path from $T$ to $Y$ exists, and treating $T$ as
observed is the same as treating it as intervened upon.

In observational studies, no such physical severance occurs.  The
treatment was not assigned by the researcher; it was selected --- by
patients, firms, governments, or individuals --- based on their
characteristics.  Those characteristics may include unmeasured variables
$U$ that also affect the outcome, opening back-door paths that
observational data alone cannot close.

The key insight of this chapter is that \emph{both settings target the
same estimand} --- the ATE $\E[Y(1) - Y(0)]$ --- but reach it by
fundamentally different routes:

\begin{tcolorbox}[defstyle, title={Two Routes to Ignorability}]
\renewcommand{\arraystretch}{1.6}
\begin{tabular}{p{3.0cm} p{4.6cm} p{4.6cm}}
\toprule
  & \textbf{Randomized experiment}
  & \textbf{Observational study} \\
\midrule
How ignorability arises
  & By design: the researcher severs all arrows into $T$
  & By assumption: the analyst asserts no unmeasured confounders \\
Status of $(Y(0),Y(1)) \indep T \mid X$
  & \emph{Guaranteed} --- holds unconditionally ($X$ not needed)
  & \emph{Assumed} --- requires $X$ to capture every confounder \\
Credibility
  & As strong as the randomization protocol
  & As strong as substantive knowledge of the data-generating process \\
Testability
  & Can audit the assignment mechanism
  & Cannot be verified from the data alone \\
Failure mode
  & Protocol violations, non-compliance
  & Any unmeasured variable that affects both $T$ and $Y$ \\
\bottomrule
\end{tabular}
\renewcommand{\arraystretch}{1}
\end{tcolorbox}

\medskip\noindent\textbf{The central pedagogical point.}  When
statisticians speak of ``controlling for confounders'' or ``adjusting
for $X$,'' they are invoking the observational route.  Statistical
adjustment can \emph{implement} ignorability once it is assumed, but it
cannot \emph{create} it.  No amount of covariate adjustment --- however
flexible the model --- can close a back-door path through an unmeasured
variable.  This is why a well-conducted randomized experiment is
considered more credible than even the most carefully adjusted
observational study: the former makes ignorability a consequence of the
design; the latter asks us to take it on faith.

\medskip\noindent The three methods developed in §§\ref{w5:sec:regression}--\ref{w5:sec:stratification}
(regression adjustment, standardization, stratification) all
\emph{implement} the observational route --- they assume ignorability
and then estimate the identified quantity.  Chapter~6 extends this with
the propensity score.  Chapter~7 abandons the ignorability route
entirely when unobserved confounders are present and turns instead to
instrumental variables.

\subsection{Definition}

With this contrast in mind, the formal definition captures the
observational case: what must hold about the treatment assignment
mechanism for causal identification to succeed without randomization?

\begin{defn}{Ignorability (Unconfoundedness)}
\label{w5:defn:ignorability}
Let $X$ denote a vector of pre-treatment covariates.  We say that
treatment assignment is \textbf{weakly ignorable} given $X$ if:
\[
  Y(t) \;\indep\; T \;\Big|\; X, \qquad t \in \{0,1\}.
\]
We say assignment is \textbf{strongly ignorable} given $X$
(Rosenbaum \& Rubin, 1983) if, in addition, the overlap condition holds:
\[
  0 < P(T{=}1 \mid X) < 1 \quad \text{almost surely.}
\]
\end{defn}

\noindent Ignorability is also called \emph{unconfoundedness},
\emph{selection on observables}, \emph{no unmeasured confounders}, or
\emph{conditional exchangeability}.  All of these phrases refer to the
same condition: conditional on $X$, the assignment mechanism is
ignorable for the purpose of causal inference.  In a randomized
experiment, weak ignorability holds \emph{unconditionally} ---
$X$ is not needed, and the condition reduces to $(Y(0), Y(1)) \indep T$,
the strongest possible form.

\subsection{Graphical Interpretation}

In the DAG language, ignorability corresponds exactly to the
\emph{back-door criterion} (Definition~\ref{w3:def:backdoor},
Theorem~\ref{w3:thm:backdoor}).  A set $X$ satisfies the back-door
criterion for the effect of $T$ on $Y$ if:
\begin{enumerate}
  \item $X$ blocks every back-door path from $T$ to $Y$ (every path
    that begins with an arrow \emph{into} $T$), and
  \item $X$ contains no descendant of $T$.
\end{enumerate}
When $X$ satisfies the back-door criterion,
$Y(t) \indep T \mid X$ holds and the back-door adjustment formula
identifies the causal effect (Theorem~\ref{w3:thm:backdoor}).

Chapter~1 introduced the \emph{causal trinity}: SEM, DAGs, and
potential outcomes as three languages for expressing the same causal
content (Figure~\ref{w1:fig:trinity}).  Ignorability has a precise
translation in all three.

\begin{tcolorbox}[defstyle, title={Three Languages for Ignorability}]
\renewcommand{\arraystretch}{1.55}
\begin{tabular}{p{3.2cm} p{9.2cm}}
\toprule
\textbf{Language} & \textbf{Statement of ignorability} \\
\midrule
SEM
  & The structural error $\varepsilon_Y$ in the outcome equation
    $Y = f_Y(T, X, \varepsilon_Y)$ is independent of $T$ given $X$:
    no unobserved variable enters both the equation that determines $T$
    and the equation that determines $Y$.
    Formally: $\varepsilon_Y \indep T \mid X$
    (cf.\ Chapter~1, Definition of SEM). \\
DAG / do-calculus
  & $X$ satisfies the back-door criterion for the effect of $T$ on $Y$:
    $X$ blocks all back-door paths from $T$ to $Y$ and contains no
    descendant of $T$
    (Definition~\ref{w3:def:backdoor}, Theorem~\ref{w3:thm:backdoor}). \\
Potential outcomes
  & $Y(t) \indep T \mid X$ for $t \in \{0,1\}$
    (weak ignorability; Definition~\ref{w5:defn:ignorability}). \\
\bottomrule
\end{tabular}
\renewcommand{\arraystretch}{1}
\end{tcolorbox}

\noindent The three statements are equivalent under the standard
bridge assumptions between frameworks (consistency + no interference;
see Chapter~1, Theorem~\ref{w1:prop:bridge}).  In practice, which
language is most useful depends on the task: the SEM formulation is
natural when writing down a structural model and asking whether any
omitted variable enters both equations; the DAG formulation is natural
for reading off the criterion from a graph; the potential outcomes
formulation is natural for writing down estimators and proofs.

\subsection{What Ignorability Requires}

Ignorability is a \emph{substantive}, not a statistical, assumption.
It requires that the analyst has measured every variable that jointly
influences treatment receipt and the outcome.  This is often called
``no unmeasured confounders.''

\begin{warn}{Ignorability is Untestable from Observational Data}
No statistical test can confirm or refute unconfoundedness using
observed data alone.  If an unobserved variable $U$ affects both $T$
and $Y$, the back-door path $T \leftarrow U \rightarrow Y$ is open, and
conditioning on any set of observed variables cannot close it.
Sensitivity analysis (Rosenbaum, 2002) can quantify how large such an
unobserved confounder would have to be to overturn a conclusion, but
cannot confirm the assumption itself.
\end{warn}

\subsection{From Ignorability to Identification}
\label{w5:subsec:identification}

Chapter~4 stated the back-door adjustment formula as
equation~\eqref{w4:eq:ate-id},
deferring the proof because it requires the law of iterated expectations
together with consistency, which were introduced only informally there.
We now supply the complete argument.  It is not merely asserted: it
follows from the three components of strong ignorability --- weak
ignorability, overlap, and the consistency assumption --- in a way that
makes explicit which step invokes which assumption and how each can
fail.  We state the key step as a separate proposition so the logical
structure is transparent.

\begin{tcolorbox}[defstyle, title={Consistency Assumption}]
\label{w5:assump:sutva}
The \textbf{consistency assumption} (part of SUTVA) states that the
observed outcome equals the potential outcome corresponding to the
treatment actually received:
\[
  Y_i = Y_i(T_i), \qquad \text{i.e.,}\quad
  T_i = t \;\Rightarrow\; Y_i = Y_i(t).
\]
It rules out interference between units and hidden versions of
treatment, and is required whenever we move between the potential
outcomes world and the observed data world.
\end{tcolorbox}

\begin{lemma}{Identification of $\E[Y(t)]$}
\label{w5:prop:ident-pot}
Suppose:
\begin{enumerate}
  \item[(i)] \emph{Weak ignorability}: $Y(t) \indep T \mid X$;
  \item[(ii)] \emph{Overlap}: $P(T{=}t \mid X{=}x) > 0$
    for $p$-almost every $x$;
  \item[(iii)] \emph{Consistency}: $Y = Y(T)$.
\end{enumerate}
Then:
\begin{equation}
  \E\bigl[Y(t)\bigr]
  = \E\Bigl[\E\bigl[Y \mid T{=}t,\, X\bigr]\Bigr]
  = \int \E\bigl[Y \mid T{=}t,\, X{=}x\bigr]\, p(x)\, dx.
  \label{w5:eq:ident-pot}
\end{equation}
\end{lemma}

\begin{proof}
Apply the law of iterated expectations, then use each assumption in
turn:
\begin{align}
  \E\bigl[Y(t)\bigr]
  &= \E\Bigl[\E\bigl[Y(t) \mid X\bigr]\Bigr]
  \tag{LIE} \\[4pt]
  &= \E\Bigl[\E\bigl[Y(t) \mid T{=}t,\, X\bigr]\Bigr]
  \tag{weak ignorability (i)} \\[4pt]
  &= \E\Bigl[\E\bigl[Y \mid T{=}t,\, X\bigr]\Bigr].
  \tag{consistency (iii)}
\end{align}
The step labelled ``weak ignorability'' uses $Y(t) \indep T \mid X$,
which gives $\E[Y(t) \mid X] = \E[Y(t) \mid T{=}t, X]$ for any value
of $T$ that has positive probability given $X$ --- and overlap
assumption~(ii) guarantees that $P(T{=}t \mid X{=}x) > 0$ almost
surely, so the conditional expectation $\E[Y \mid T{=}t, X{=}x]$ is
well-defined everywhere in the support of $X$.  The step labelled
``consistency'' uses $Y = Y(t)$ on the event $\{T = t\}$ to replace
$\E[Y(t) \mid T{=}t, X]$ by $\E[Y \mid T{=}t, X]$.
\end{proof}

\begin{remark}[Two Languages, One Formula]
\label{w5:rem:two-languages}
Lemma~\ref{w5:prop:ident-pot} and Theorem~\ref{w3:thm:backdoor}
are two proofs of the same identification formula, written in different
languages.

\smallskip
\noindent\textit{Theorem~\ref{w3:thm:backdoor} (Chapter~3)} works in
the graphical language.  The assumption is that the adjustment set
$\mathbf{S}$ satisfies the back-door criterion in the DAG $\Gcal$:
every back-door path from $T$ to $Y$ is blocked.  The conclusion is
$P(y \mid \doop(T{=}t)) = \int P(y \mid T{=}t, \mathbf{S}{=}\mathbf{s})
\, dP(\mathbf{s})$, proved by applying Rules~2 and~3 of the
do-calculus in the mutilated graph $\Gcal_{\overline{T}}$.

\smallskip
\noindent\textit{Lemma~\ref{w5:prop:ident-pot} (this section)}
works in the potential-outcomes language.  The assumption is strong
ignorability: $Y(t) \indep T \mid X$ (weak ignorability), $P(T{=}t
\mid X) > 0$ almost surely (overlap), and $Y = Y(T)$ (consistency).
The conclusion is $\E[Y(t)] = \E[\E[Y \mid T{=}t, X]]$, proved via the
law of iterated expectations.

\smallskip
\noindent\textit{Why they agree.}  The bridge is the consistency
assumption: $Y_i = Y_i(T_i)$ implies $\E[Y(t)] = \E[Y \mid
\doop(T{=}t)]$, so the potential-outcomes target and the
interventional target are the same quantity.  On the graphical side,
when $\mathbf{S} = X$ satisfies the back-door criterion, it can be
shown that $Y(t) \indep T \mid X$ holds in the corresponding
non-parametric structural equation model --- this is the formal
equivalence between graphical and potential-outcomes conditions
established by \citet{pearl2009causality}.  Chapter~3 proves the formula
from the graph; this chapter proves it from the potential outcomes.
Practitioners trained in one tradition will recognize the same result
stated in the other.
\end{remark}

\noindent The ATE and ATT follow immediately as corollaries.

\begin{thm}{Identification of the ATE and ATT}
\label{w5:prop:ident-ate}
Under the same three assumptions:
\begin{align}
  \tau_{\mathrm{ATE}}
  &= \E\bigl[Y(1) - Y(0)\bigr]
  = \int \bigl[\E(Y \mid T{=}1, X{=}x)
             - \E(Y \mid T{=}0, X{=}x)\bigr]\, p(x)\, dx,
  \label{w5:eq:backdoor} \\[6pt]
  \tau_{\mathrm{ATT}}
  &= \E\bigl[Y(1) - Y(0) \mid T{=}1\bigr]
  = \int \bigl[\E(Y \mid T{=}1, X{=}x)
             - \E(Y \mid T{=}0, X{=}x)\bigr]\, p(x \mid T{=}1)\, dx.
  \label{w5:eq:backdoor-att}
\end{align}
\end{thm}

\begin{proof}
Apply Lemma~\ref{w5:prop:ident-pot} separately to $t = 1$ and
$t = 0$ and subtract.  For the ATT, restrict to the subpopulation
$T = 1$: replace $p(x)$ by the conditional density $p(x \mid T{=}1)$
in the same argument, noting that overlap for $t = 0$ still requires
$P(T{=}0 \mid X{=}x) > 0$ on the support of $X \mid T{=}1$.
\end{proof}

\medskip\noindent\textbf{What the proof uses and where it can fail.}
Table~\ref{w5:tab:ident-steps} traces each step in the proof back to
the assumption it invokes, and gives the failure mode if the assumption
is violated.

\begin{table}[h]
\centering
\renewcommand{\arraystretch}{1.5}
\begin{tabular}{p{3.5cm} p{4.0cm} p{4.8cm}}
\toprule
\textbf{Proof step} & \textbf{Assumption invoked} & \textbf{Failure mode} \\
\midrule
$\E[Y(t) \mid X] = \E[Y(t) \mid T{=}t, X]$
  & Weak ignorability (i)
  & Unmeasured confounder $U$: $Y(t)$ depends on $T$ even within $X$-strata \\
$\E[Y \mid T{=}t, X{=}x]$ is well-defined
  & Overlap (ii)
  & Empty stratum: no units with $T{=}t$ at covariate value $x$; estimator extrapolates \\
$\E[Y(t) \mid T{=}t, X] = \E[Y \mid T{=}t, X]$
  & Consistency (iii)
  & Interference or hidden treatment versions: $Y_i$ depends on $T_j$ for $j \ne i$ \\
\bottomrule
\end{tabular}
\caption{The three assumptions underlying identification of $\E[Y(t)]$,
and what goes wrong when each is violated.}
\label{w5:tab:ident-steps}
\end{table}

\begin{remark}
Overlap is \emph{testable} from data: for any candidate $X$, one can
check whether the distributions of $X \mid T{=}1$ and $X \mid T{=}0$
have common support.  Weak ignorability and consistency are not
testable from data alone --- they are substantive assumptions about the
data-generating process.  This asymmetry is important: overlap
violations can be diagnosed and addressed (e.g.\ by restricting the
estimand to the region of common support); ignorability violations
cannot be detected and must be justified by the study design and domain
knowledge.  A full treatment of overlap diagnostics using the
propensity score is given in Chapter~6.
\end{remark}

\noindent Equations \eqref{w5:eq:backdoor}--\eqref{w5:eq:backdoor-att}
are the \emph{identification results}.  The next three sections present
methods for \emph{estimating} the identified quantities from data.

\subsection{Which Variables to Condition On: The Pre-Treatment Requirement}
\label{w5:subsec:pretreatment}

The back-door formula conditions on $X$.  Before estimating anything,
there is a prior question that is purely graphical: \emph{which
variables should be in $X$?}  The answer is not ``all available
variables.''  Conditioning on the wrong variable can introduce bias
rather than remove it.  The back-door criterion gives the correct
answer, and it imposes one requirement that is easy to state but easy
to forget in practice.

\begin{warn}{Never Condition on a Post-Treatment Variable}
\label{w5:warn:posttreatment}
A variable $L$ is \emph{post-treatment} if it is caused by $T$, i.e.\
$T \to L$ in the DAG.  Conditioning on $L$ in the outcome regression
--- including $L$ in $X$ --- can bias the estimated treatment effect,
even if $L$ would otherwise be a strong predictor of $Y$.  There are
two distinct mechanisms:

\begin{enumerate}
  \item \textbf{Blocking a causal pathway.}  If $L$ lies on a directed
    path from $T$ to $Y$ (i.e.\ $L$ is a mediator: $T \to L \to Y$),
    conditioning on $L$ blocks part of the causal effect of $T$.  The
    back-door formula applied to $(T, Y)$ with $L$ in $X$ estimates a
    \emph{direct} effect, not the total effect.  Whether this is the
    right estimand depends on the research question; mistaking a direct
    effect for a total effect is a common error.

  \item \textbf{Opening a collider path.}  If $L$ has two parents ---
    $T \to L \leftarrow U$ --- then $L$ is a \emph{collider} on the
    path $T \to L \leftarrow U \to Y$.  This path is blocked when $L$
    is not conditioned on.  Conditioning on $L$ \emph{opens} the path,
    creating a spurious association between $T$ and $Y$ through $U$
    that did not exist before adjustment.
\end{enumerate}

In both cases, including $L$ in $X$ violates condition (i) of the
back-door criterion: \emph{$X$ must contain no descendant of $T$.}
The back-door criterion is therefore not just a sufficiency condition
for identification --- it is also the correct filter for deciding
which variables to include.
\end{warn}

\begin{figure}[h]
\centering
\begin{tikzpicture}[
  node distance=14mm and 20mm,
  nd/.style={circle, draw=accent, line width=0.9pt,
             minimum size=9mm, font=\small, fill=defbg},
  un/.style={circle, draw=backdoorred, line width=0.9pt, dashed,
             minimum size=9mm, font=\small, fill=warnbg},
  ar/.style={-{Stealth[length=4pt]}, line width=0.9pt, color=accent},
  ra/.style={-{Stealth[length=4pt]}, line width=0.9pt, color=backdoorred, dashed},
  lbl/.style={font=\footnotesize\itshape\color{darkgrey}}
]
%% ── Left: Mediator ────────────────────────────────────────────────
\begin{scope}[xshift=0cm]
  \node[nd] (T1) {$T$};
  \node[nd, right=20mm of T1] (L1) {$L$};
  \node[nd, right=20mm of L1] (Y1) {$Y$};
  \draw[ar] (T1) -- (L1);
  \draw[ar] (L1) -- (Y1);
  \draw[ar] (T1) to[bend right=30] (Y1);
  \node[lbl, below=5mm of L1] {mediator};
  \node[lbl, below=14mm of T1, xshift=22mm, align=center]
    {\textbf{Mediator}: $T \to L \to Y$\\
     conditioning on $L$ blocks\\
     the indirect pathway};
\end{scope}
%% ── Right: Collider ───────────────────────────────────────────────
\begin{scope}[xshift=7.5cm]
  \node[nd] (T2) {$T$};
  \node[nd, right=20mm of T2] (L2) {$L$};
  \node[nd, right=20mm of L2] (Y2) {$Y$};
  \node[un, above=14mm of L2] (U2) {$U$};
  \draw[ar] (T2) -- (L2);
  \draw[ra] (U2) -- (L2);
  \draw[ra] (U2) -- (Y2);
  \node[lbl, below=5mm of L2] {collider};
  \node[lbl, below=14mm of T2, xshift=22mm, align=center]
    {\textbf{Collider}: $T \to L \leftarrow U \to Y$\\
     conditioning on $L$ opens\\
     a spurious $T$--$Y$ path};
\end{scope}
\end{tikzpicture}
\caption{Two ways post-treatment conditioning introduces bias.
\textbf{Left}: $L$ is a mediator ($T \to L \to Y$); conditioning on
$L$ blocks the indirect pathway and yields a direct effect rather than
the total effect.  \textbf{Right}: $L$ is a collider on the path
$T \to L \leftarrow U \to Y$; conditioning on $L$ opens a previously
blocked path and creates spurious $T$--$Y$ association through the
unobserved $U$.  In both cases $L$ is a descendant of $T$, violating
the back-door criterion.}
\label{w5:fig:posttreatment}
\end{figure}

\medskip\noindent\textbf{The practical rule.}  Before running any
regression or computing any cell means, classify every candidate
covariate as pre-treatment or post-treatment using the causal graph.
Only pre-treatment variables that satisfy the back-door criterion
belong in $X$.  Post-treatment variables should be left out unless the
research question is specifically about a direct or mediated effect ---
in which case Chapter~8 provides the appropriate framework.

% ============================================================
\section{Regression Adjustment and Standardization}
\label{w5:sec:regression}
% ============================================================

\subsection{The Common Three-Step Logic}

Both regression adjustment and standardization are implementations of
the same back-door formula~\eqref{w5:eq:backdoor}.  They differ only in
\emph{how} they estimate the conditional mean
$\mu(t, x) = \E[Y \mid T{=}t, X{=}x]$, not in the underlying procedure.
Understanding them as two variants of the same recipe is the key to
using either one fluently.  Throughout this section $X$ denotes a set
of \emph{pre-treatment} covariates satisfying the back-door criterion;
the preceding section explains why post-treatment variables must be
excluded.

\begin{tcolorbox}[defstyle, title={The Three-Step Recipe: Outcome Regression / G-Formula}]
\begin{enumerate}
  \item \textbf{Estimate the outcome model.}
    Fit a model for $\mu(t,x) = \E[Y \mid T{=}t, X{=}x]$ using the
    observed data $\{(Y_i, T_i, X_i)\}_{i=1}^n$.  Any regression
    method may be used: OLS, logistic regression, a flexible
    nonparametric estimator, or cell means (see §\ref{w5:subsec:std}).

  \item \textbf{Predict both potential outcomes for every unit.}
    For each unit $i$ in the sample --- \emph{regardless of the treatment
    that unit actually received} --- compute:
    \[
      \hat{Y}_i(1) = \hat\mu(1, X_i),
      \qquad
      \hat{Y}_i(0) = \hat\mu(0, X_i).
    \]
    This is the counterfactual imputation step: for a treated unit,
    $\hat{Y}_i(0)$ is the predicted outcome had that unit been assigned
    to control; for a control unit, $\hat{Y}_i(1)$ is the predicted
    outcome had that unit been assigned to treatment.

  \item \textbf{Average the individual treatment effect estimates.}
    The ATE estimator is the sample average of the imputed individual
    differences:
    \begin{equation}
      \hat\tau_{\mathrm{OR}}
      = \frac{1}{n}\sum_{i=1}^n
        \bigl[\hat{Y}_i(1) - \hat{Y}_i(0)\bigr]
      = \frac{1}{n}\sum_{i=1}^n
        \bigl[\hat\mu(1, X_i) - \hat\mu(0, X_i)\bigr].
      \label{w5:eq:or}
    \end{equation}
    To estimate the ATT instead, average only over the $n_1$ treated
    units: $\hat\tau_{\mathrm{ATT}} = \frac{1}{n_1}\sum_{i:\,T_i=1}
    [\hat\mu(1,X_i) - \hat\mu(0,X_i)]$.
\end{enumerate}
\end{tcolorbox}

\noindent This estimator is called the \emph{outcome regression} (OR)
estimator or the \emph{G-computation} estimator (Robins, 1986).  The
G-formula name emphasizes that step~3 \emph{averages out} the
covariates $X$ by plugging in the empirical distribution of $X$,
exactly as the back-door formula~\eqref{w5:eq:backdoor} requires.

\medskip\noindent\textbf{Why the three steps correspond to the
identifying formula.}  The back-door formula says:
\[
  \tau_{\mathrm{ATE}}
  = \int \bigl[\mu(1,x) - \mu(0,x)\bigr]\, p(x)\, dx.
\]
Step~1 estimates $\mu$; step~2 evaluates $\mu(1,X_i) - \mu(0,X_i)$
at the observed covariate value of each unit; step~3 approximates the
integral by the sample average, using the empirical distribution of
$X$ as a nonparametric estimate of $p(x)$.  The estimator is
\emph{plug-in}: replace every population quantity in the identification
formula with its sample counterpart.

\subsection{A Worked Example}

Consider a binary covariate $X \in \{0,1\}$ (e.g.\ sex) and a
continuous outcome $Y$ (e.g.\ earnings in thousands of dollars).  An
OLS model fit separately within each treatment arm gives:

\medskip
\begin{center}
\renewcommand{\arraystretch}{1.3}
\begin{tabular}{lcccc}
\toprule
& \multicolumn{2}{c}{\textbf{Treated arm} ($T=1$)}
& \multicolumn{2}{c}{\textbf{Control arm} ($T=0$)} \\
$X$ & $n$ & $\bar{Y}$ & $n$ & $\bar{Y}$ \\
\midrule
$0$ & 30 & 42 & 70 & 35 \\
$1$ & 70 & 58 & 30 & 50 \\
\midrule
\textbf{All} & 100 & 53.2 & 100 & 38.5 \\
\bottomrule
\end{tabular}
\end{center}
\medskip

\noindent The unadjusted difference in means is
$53.2 - 38.5 = 14.7$.  But $X{=}1$ units (with higher baseline
earnings) are overrepresented in the treated arm: 70\% treated vs.\
30\% control.  $X$ confounds the comparison.

\medskip\noindent\textbf{Step 1.}  The estimated conditional means are
$\hat\mu(1, 0) = 42$, $\hat\mu(1, 1) = 58$, $\hat\mu(0, 0) = 35$,
$\hat\mu(0, 1) = 50$.

\medskip\noindent\textbf{Step 2.}  The imputed within-stratum effects
are: $\hat\mu(1,0) - \hat\mu(0,0) = 42 - 35 = 7$ (for $X{=}0$ units)
and $\hat\mu(1,1) - \hat\mu(0,1) = 58 - 50 = 8$ (for $X{=}1$ units).

\medskip\noindent\textbf{Step 3.}  Averaging over the \emph{marginal}
distribution of $X$: 50\% of the full sample has $X{=}0$, 50\% has
$X{=}1$ (equal mix in the combined sample of 200), giving:
\[
  \hat\tau_{\mathrm{ATE}}
  = 7 \times 0.5 + 8 \times 0.5 = 7.5.
\]
After adjusting for $X$, the estimated treatment effect is 7.5, not
14.7.  The difference arises because the unadjusted comparison
conflates the treatment effect with the higher baseline earnings of
$X{=}1$ units who happen to be treated more often.

\subsection{Standardization as a Special Case}
\label{w5:subsec:std}

Standardization is the name traditionally used in epidemiology and
biostatistics for the special case of regression adjustment in which
$X$ is categorical with finite support $\mathcal{X}$.  When $X$ is
categorical, no parametric model for $\mu(t,x)$ is needed: the
conditional mean $\hat\mu(t,x)$ is simply the sample mean of $Y$
within each cell $(T{=}t, X{=}x)$.  This is equivalent to fitting a
fully saturated outcome model.

The resulting estimator is:
\begin{equation}
  \hat\tau_{\mathrm{STD}}
  = \sum_{x \in \mathcal{X}}
    \bigl[\hat\mu(1,x) - \hat\mu(0,x)\bigr]\,\hat{p}(x),
  \label{w5:eq:std}
\end{equation}
where $\hat\mu(t,x) = \bar{Y}_{t,x}$ is the sample mean of $Y$ among
units with $T{=}t$ and $X{=}x$, and $\hat{p}(x) = n_x/n$ is the
empirical marginal proportion.  Equation~\eqref{w5:eq:std} is
\emph{algebraically identical} to the plug-in form of
\eqref{w5:eq:or}: both compute $\frac{1}{n}\sum_i [\hat\mu(1,X_i) -
\hat\mu(0,X_i)]$.  Standardization is therefore regression adjustment
with a saturated outcome model, and the two estimators coincide
whenever $X$ is categorical.

The worked example above \emph{is} a standardization: the cell means
$\hat\mu(t,x)$ were computed directly from the data without any
parametric model.  The formula is also known as the \emph{G-formula}
(Robins, 1986) and as \emph{direct standardization} in epidemiology.

\begin{tcolorbox}[defstyle, title={Regression Adjustment vs.\ Standardization: Same Procedure, Different Estimator of $\hat\mu$}]
\renewcommand{\arraystretch}{1.4}
\begin{tabular}{p{3.5cm} p{4.5cm} p{4.5cm}}
\toprule
& \textbf{Regression adjustment}
& \textbf{Standardization} \\
\midrule
How $\hat\mu(t,x)$ is estimated
  & Parametric model: OLS, logistic, or flexible learner
  & Cell means: $\bar Y$ within $(T{=}t, X{=}x)$ \\
Works when $X$ is
  & Continuous or high-dimensional
  & Discrete and low-dimensional \\
Bias if wrong
  & Model misspecification
  & Sparse cells (some $(t,x)$ strata empty) \\
Steps 2--3
  & Identical: predict both potential outcomes, average
  & Identical: predict both potential outcomes, average \\
\bottomrule
\end{tabular}
\renewcommand{\arraystretch}{1}
\end{tcolorbox}

\noindent When $X$ is continuous, cell means are unavailable.
$\hat\mu(t,x)$ is then estimated by a regression model --- OLS,
logistic, or a nonparametric learner --- and the three-step recipe
applies with no other change.

\subsection{Model Specification and What Can Go Wrong}

The OR estimator is consistent if and only if $\hat\mu(t,x) \to
\mu(t,x)$ as $n \to \infty$ --- that is, if the outcome model is
correctly specified.  This is the main vulnerability.

\medskip\noindent\textbf{Misspecification bias.}  If the true
$\mu(t,x)$ is nonlinear but a linear model is used, the estimated
treatment effect absorbs part of the functional form error.  The bias
can be in either direction.  Flexible machine learning estimators for
$\hat\mu$ reduce this risk, at the cost of requiring cross-fitting
(sample splitting) to avoid overfitting bias in the second step; this
is taken up in Chapter~11.

\medskip\noindent\textbf{Extrapolation.}  Step~2 predicts
$\hat{Y}_i(0)$ for treated units with covariate values $X_i$ that may
lie outside the support of the control group.  In regions with no
control data, $\hat\mu(0, X_i)$ is extrapolated from the model rather
than estimated from comparable units.  The propensity score overlap
condition (Chapter~6) formalises when extrapolation is unavoidable.

\subsection{Regression Adjustment in Randomized Experiments: Lin (2013)}

In a \emph{randomized} experiment, outcome regression can be used to
reduce variance even though covariate adjustment is not needed for
unbiasedness.  The \emph{regression-adjusted estimator} of
\citet{lin2013agnostic} fits the fully interacted OLS model:
\[
  Y_i = \alpha + \beta T_i + \gamma^\top \tilde{X}_i
        + \delta^\top (T_i \cdot \tilde{X}_i) + \varepsilon_i,
\]
where $\tilde{X}_i = X_i - \bar{X}$ are mean-centered covariates.  The
OLS coefficient $\hat\beta$ is read off as the ATE estimator.

\begin{thm}{Design-Consistency of the Regression-Adjusted Estimator \citep{lin2013agnostic}}
\label{w5:prop:lin}
Under complete randomization, $\hat\beta$ from the interacted
regression is consistent for $\tau_{\mathrm{ATE}}$ regardless of
whether the linear model is correctly specified.
\end{thm}

\noindent\textit{Sketch.}  Because $T \indep X$ under randomization,
the projection of $Y$ on $(T, X, T \cdot \tilde{X})$ recovers the
treatment coefficient without confounding bias even if the true $\mu$
is not linear.  The interaction terms ensure that the estimated
treatment effect is the sample analog of the semiparametric efficient
estimator.  For the full proof see \citet{lin2013agnostic} or
\citet{freedman2008regression}.

Importantly, $\hat\beta$ from the Lin estimator is numerically
equal to the three-step OR estimator~\eqref{w5:eq:or} applied with a
linear model fit separately within each treatment arm.  The Lin
estimator is therefore not a separate method: it is regression
adjustment, presented in a single-regression form that makes inference
straightforward.

\begin{remark}[The Role of the Regression Model in a Randomized Experiment]
\label{w5:rem:reg-role}
In a randomized experiment the regression model serves one purpose
only: \emph{variance reduction}.  It does not contribute to
unbiasedness.

Unbiasedness (or, in large samples, design-consistency) is guaranteed
by the randomization alone, regardless of whether the outcome model is
correctly specified --- this is precisely what Theorem~\ref{w5:prop:lin}
states.  The plain difference-in-means estimator $\hat\tau_{\mathrm{DIM}}$
is already design-consistent; the only reason to add covariates is to
reduce $\mathrm{Var}(\hat\tau)$ by absorbing residual variation in $Y$
that is unrelated to $T$.  If the regression model fits well, the
residual variance shrinks and the estimator becomes more precise.  If
the model is misspecified or the covariates are weakly predictive, the
variance reduction is small or zero, but no bias is introduced.

This stands in sharp contrast to the observational setting
(Section~\ref{w5:sec:ignorability}), where the regression model
carries a \emph{double} burden: it must both adjust for confounding
(unbiasedness) and fit the outcome surface well (efficiency).
Misspecification in an observational study therefore threatens
\emph{both} properties simultaneously.  In a randomized experiment, it
threatens only efficiency.
\end{remark}

\begin{remark}[Precedence in Survey Sampling]
\label{w5:rem:lin-survey}
The design-consistency result in Theorem~\ref{w5:prop:lin} --- that the
regression-adjusted estimator converges in probability to
$\tau_{\mathrm{ATE}}$ under the randomization distribution regardless
of whether the working linear model is correctly specified --- was
established in the survey sampling literature decades before
\citet{lin2013agnostic}.

\citet{isaki1982survey} introduced the \emph{generalized regression
(GREG) estimator} for finite-population means under general probability
sampling designs, including simple random sampling, and proved that it
is design-consistent under the randomization distribution even when the
assisting regression model is misspecified.  Because a completely
randomized experiment is structurally equivalent to simple random
sampling from the finite population $\mathcal{F}_N$ of potential
outcomes (see Section~\ref{w5:sec:rct}), their result directly implies
Theorem~\ref{w5:prop:lin}.

\citet{deville1992calibration} extended this class further to
\emph{calibration estimators}, which adjust the sampling weights so
that weighted sample totals match known population totals for the
auxiliary variables $X$.  The regression estimator is the special case
of a calibration estimator using linear distance; all calibration
estimators share the design-consistency property of Isaki and Fuller.

\citet{lin2013agnostic}'s contribution was to restate and prove this
result within the Neyman--Rubin potential-outcomes framework, connect
it to semiparametric efficiency theory, and show that the
Huber--White sandwich variance estimator yields asymptotically valid
confidence intervals under the randomization distribution --- a
presentation tailored to the experimental causal inference audience
where these earlier survey sampling results were not widely known.
\end{remark}

\begin{remark}
Lin's consistency result is specific to randomized experiments.  In an
observational study, $\hat\tau_{\mathrm{OR}}$ is consistent only if
$\hat\mu$ converges to the true $\mu$.  The doubly-robust (AIPW)
estimator of Chapter~10 achieves consistency if \emph{either} the
outcome model or the propensity score model is correctly specified,
providing a robustness guarantee the plain OR estimator lacks.
\end{remark}

% ============================================================
\section{Stratification}
\label{w5:sec:stratification}
% ============================================================

\subsection{The Idea}

Stratification (subclassification) is a non-parametric implementation
of the back-door formula.  Rather than modeling $\E[Y \mid T, X]$, the
analyst divides the sample into \emph{strata} --- subgroups of units
with the same or similar values of $X$ --- and estimates the treatment
effect within each stratum.

Within a stratum $\mathcal{S}_k = \{i : X_i \in B_k\}$, the treated and
control units have similar covariate distributions.  If the stratum is
narrow enough (i.e.\ $B_k$ is small), ignorability holds approximately
within the stratum, and the within-stratum difference in means:
\[
  \hat\tau_k
  = \bar Y_{1,k} - \bar Y_{0,k}
  = \frac{\sum_{i \in \mathcal{S}_k} T_i Y_i}{\sum_{i \in \mathcal{S}_k} T_i}
  - \frac{\sum_{i \in \mathcal{S}_k} (1-T_i) Y_i}{\sum_{i \in \mathcal{S}_k} (1-T_i)}
\]
is an approximately unbiased estimate of the stratum-specific ATE.

\subsection{The Stratified Estimator}

The overall ATE is estimated as a weighted average of within-stratum
estimates:
\begin{equation}
  \hat\tau_{\mathrm{STRAT}}
  = \sum_{k=1}^K \hat\tau_k \cdot \frac{n_k}{n},
  \label{w5:eq:strat}
\end{equation}
where $K$ is the number of strata, $n_k$ is the number of units in
stratum $k$, and $n = \sum_k n_k$.  Cochran (1968) showed that with
$K = 5$ strata of equal size, roughly 90\% of the bias from a single
continuous confounder is removed; with $K = 10$ strata, over 95\%.

\begin{remark}[Stratification, Standardization, and Survey Sampling Terminology]
\label{w5:rem:strat-poststrat}
When $X$ is categorical, the stratified estimator~\eqref{w5:eq:strat}
and the standardization estimator~\eqref{w5:eq:std} are algebraically
identical: both compute within-cell treatment effect estimates and
aggregate them with marginal weights $n_k/n$.  The two names reflect
disciplinary tradition rather than a methodological distinction.

When $X$ is continuous, stratification estimates $\hat\mu(t,x)$ by a
piecewise-constant step function (one constant per stratum), while
regression adjustment fits a smooth parametric or nonparametric model.
As the number of strata grows and stratum widths shrink, the stratified
estimator converges to the regression-adjusted estimator.

In the survey sampling literature, \emph{stratification} and
\emph{post-stratification} are distinct operations.
\emph{Stratification} is a \emph{design-stage} decision: the
population is divided into strata before sampling, and units are
sampled independently within each stratum to control the allocation.
\emph{Post-stratification} is an \emph{analysis-stage} adjustment:
stratum membership is recorded after a non-stratified sample has been
drawn, and the estimator is reweighted by marginal stratum counts ---
exactly what equation~\eqref{w5:eq:std} does.  Standardization in
causal inference corresponds to post-stratification, not to
stratified sampling design.
\end{remark}

\subsection{Limitations}

Stratification faces the \emph{curse of dimensionality}: with $p$
binary covariates there are $2^p$ cells, many of which will be empty in
practice.  This motivates two solutions studied in this course:
\begin{itemize}
  \item \textbf{Regression adjustment / standardization} (§\ref{w5:sec:regression}) imposes
    parametric structure on $\hat\mu(t,x)$, allowing interpolation
    across cells.
  \item \textbf{Propensity score stratification} (Chapter~6) collapses
    the multivariate $X$ to a scalar $\pi(X) = P(T{=}1 \mid X)$ and
    stratifies on the scalar.  Theorem~\ref{w6:thm:ps} (Chapter~6)
    guarantees that strata defined by $\pi(X)$ balance the full
    multivariate $X$.
\end{itemize}

% ============================================================
\section{Simpson's Paradox}
\label{w5:sec:simpson}
% ============================================================

Simpson's paradox --- the reversal of an association when conditioning
on a third variable --- was introduced in Chapter~1
(Section~\ref{w1:sec:two-step}, in particular
Section~\ref{w1:subsec:simpson}).  There, a flu-vaccination example
showed that the vaccine appeared harmful in the aggregate yet beneficial
within every age stratum, because elderly individuals were
disproportionately vaccinated.  Conditioning on age via the back-door
formula recovered the correct causal effect.

Chapter~5's development gives that example a second layer of
interpretation.  The back-door formula used in Chapter~1 is precisely
the standardization estimator~\eqref{w5:eq:std}: within-stratum
conditional means averaged over the \emph{marginal} distribution of the
confounder rather than its treatment-conditional distribution.  The
weighted average corrects the confounding; the pooled association does
not, because it weights strata by $P(X \mid T)$ instead of $P(X)$.

\subsection{When Should You \emph{Not} Condition on $X$?}
\label{w5:subsec:simpson-collider}

Simpson's paradox also has a mirror image: conditioning on a variable
can \emph{create} a spurious association or make a true causal effect
disappear.  This occurs when $X$ is a \emph{mediator} (a descendant of
$T$ on the causal path to $Y$) or a \emph{collider}.  The back-door
criterion gives the correct prescription: include $X$ in the adjustment
set only if it blocks back-door paths \emph{without} opening collider
paths.

\begin{warn}{Conditioning on a Mediator or Collider}
Adjusting for a post-treatment variable $X$ that lies on a causal
pathway $T \to X \to Y$ blocks part of the causal effect of $T$,
producing a downward-biased estimate of the total effect.  Adjusting
for a collider $X$ with $T \to X \leftarrow Y$ opens a spurious
association between $T$ and $Y$ that does not exist in the
population.  Neither mistake is detectable from the data alone; the
DAG is essential.
\end{warn}

\noindent We will revisit the mediator case in detail when we study
mediation analysis (Chapter~8).  The collider case was covered in
Chapter~2 and is the subject of ongoing work in sensitivity analysis.

% ============================================================
\section{Lab: Simulation Study of the Outcome Regression Estimator}
\label{w5:sec:lab}
% ============================================================

The identification results in Section~\ref{w5:sec:ignorability} and the
estimation recipes in Section~\ref{w5:sec:regression} rest on two
assumptions that pull in opposite directions: the outcome model
$\hat\mu(t,x)$ must be correctly specified for the linear outcome
regression (OR) estimator to be consistent in an observational study,
yet a more flexible nonparametric estimator pays a variance penalty for
that flexibility.  Under a completely randomized design (CRD),
Theorem~\ref{w5:prop:lin} guarantees that the linear OR estimator is
consistent regardless of model specification; the guarantee disappears
as soon as treatment is confounded.

This simulation study puts those claims side by side across two
experimental designs and two estimators, with a nonlinear
data-generating process that makes the tradeoff clearly visible.

% ----------------------------------------
\subsection*{Data-Generating Process}
\addcontentsline{toc}{subsection}{Data-Generating Process}

\begin{tcolorbox}[defstyle, title={DGP for Lab~5}]
\label{w5:lab:dgp}
Let $X \sim \mathrm{Uniform}(0,1)$ be a single pre-treatment covariate
and $\varepsilon \sim \mathcal{N}(0,1)$ independent of $X$.  The
potential outcome surfaces are:
\begin{equation}
  Y(t) \;=\; t \;+\; 3(1+t)X^2 \;+\; \varepsilon,
  \qquad t \in \{0,1\}.
  \label{w5:lab:dgp-eq}
\end{equation}
The conditional average treatment effect is $\tau(x) = 1 + 3x^2$, so
\begin{equation}
  \tau_{\mathrm{ATE}}
  = \E[1 + 3X^2]
  = 1 + 1 = 2.
  \label{w5:lab:tau}
\end{equation}
Two assignment mechanisms are studied:
\begin{itemize}
  \item \textbf{CRD:}  $T \sim \mathrm{Bern}(0.5)$, independent of $X$.
  \item \textbf{Observational:}  $T \mid X \sim
    \mathrm{Bern}\!\left(\mathrm{expit}(-2 + 5X)\right)$,
    giving $\pi(0) \approx 0.12$, $\pi(0.5) \approx 0.62$,
    $\pi(1) \approx 0.95$.  Strong confounding: high-$X$ units are
    nearly always treated.
\end{itemize}
The observed outcome is $Y = Y(T)$ in both cases.  Strong ignorability
holds in both: in the CRD because $T \indep X$; in the observational
study because treatment is a deterministic function of $X$ alone
(no unmeasured confounder).
\end{tcolorbox}

% ----------------------------------------
\subsection*{Estimators}
\addcontentsline{toc}{subsection}{Estimators}

Two outcome regression estimators of the form~\eqref{w5:eq:or} are
compared.  Both follow the three-step recipe: estimate
$\hat\mu(t,x)$, impute both potential outcomes for every unit on a
fine grid over $[0,1]$, and average.  They differ only in how
$\hat\mu(t,x)$ is estimated.

\medskip
\noindent\textbf{Linear OR.}  Within each arm $t \in \{0,1\}$, fit the
simple linear regression $Y \sim X$ by OLS and use the fitted line to
predict at all grid points.  This model is \emph{misspecified}: the
true $\mu(t,x)$ contains a quadratic term $1.5x^2$ and a
treatment--covariate interaction $tx$, neither of which is included.

\medskip
\noindent\textbf{Local-linear OR.}  Within each arm, estimate
$\mu(t,x_0)$ by local linear regression using a Gaussian kernel with
bandwidth $h = 0.5 \cdot n^{-1/5}$ (moderately undersmoothed relative
to the Silverman rule).  Local linear regression has the same
asymptotic bias order as the Nadaraya--Watson estimator but eliminates
boundary bias, which is important here because the observational
propensity score places treated and control units in different regions
of $[0,1]$.  In R, this is implemented via \texttt{KernSmooth::locpoly}.

% ----------------------------------------
\subsection*{Results}
\addcontentsline{toc}{subsection}{Results}

Table~\ref{w5:tab:lab} reports the Monte Carlo mean, bias, standard
deviation, and RMSE across $B = 2{,}000$ replications at $n = 1{,}000$.

\begin{table}[h]
\centering
\renewcommand{\arraystretch}{1.5}
\caption{Monte Carlo performance of the linear and local-linear OR
  estimators under CRD and under an observational design satisfying
  strong ignorability.  True $\tau_{\mathrm{ATE}} = 2.000$.
  $n = 1{,}000$, $B = 2{,}000$ replications, seed \texttt{2024}.}
\label{w5:tab:lab}
\begin{tabular}{llcccc}
\toprule
\textbf{Design} & \textbf{Estimator} & \textbf{Mean} & \textbf{Bias}
  & \textbf{SD} & \textbf{RMSE} \\
\midrule
CRD
  & Linear OR       & 2.0013 &  $+$0.0013 & 0.0679 & 0.0678 \\
  & Local-linear OR & 2.0333 &  $+$0.0333 & 0.0659 & 0.0738 \\
\midrule
Observational
  & Linear OR       & 1.9602 & $-$0.0398 & 0.0851 & 0.0939 \\
  & Local-linear OR & 2.0240 & $+$0.0240 & 0.0923 & 0.0953 \\
\bottomrule
\end{tabular}
\end{table}

\medskip
Three lessons stand out.

\medskip
\noindent\textbf{1.\ Under CRD, the linear OR estimator is unbiased 
regardless of model specification.}   This confirms
Theorem~\ref{w5:prop:lin}: under complete randomization, the linear
OR estimator converges in probability to $\tau_{\mathrm{ATE}}$ even
when the outcome model is wrong.  The randomization distribution,
not the model, is doing the identification work.

\medskip
\noindent\textbf{2.\ Under an observational design, the misspecified
linear OR is biased.}  The linear estimator's bias grows from $+0.0013$
under CRD to $-0.0398$ under the observational design.  The mechanism is precisely the one described in
Section~\ref{w5:sec:regression}: the omitted $X^2$ and $TX$ terms
in the outcome model cause $\hat\mu(t,x)$ to misrepresent the
outcome surface, and because high-$X$ units are predominantly treated
(the propensity score ranges from 0.12 to 0.95), the misfit is
systematically amplified in the direction of the confounding.  In a
CRD the same misfit is present, but the balanced treatment assignment
averages it out; in the observational study it does not.

\medskip
\noindent\textbf{3.\ The local-linear OR nearly eliminates bias at the
cost of higher variance.}  Under the observational design the
local-linear OR reduces bias  but
its standard deviation is 0.0923, compared to 0.0851 for the
misspecified linear model.  This variance penalty is the price of flexibility: a more expressive model uses data
less efficiently.  The RMSE comparison (0.0939 vs.\ 0.0953) therefore
favors the linear estimator in terms of mean squared error, even
though the linear estimator is biased.  This is the classic
bias--variance tradeoff, and it illustrates why a well-specified
parametric model can outperform a nonparametric one even in finite
samples.

\begin{warn}{Model Misspecification and Confounding Interact}
Comparing the two ``Linear OR'' rows: the bias is $+0.0013$ under CRD
and $-0.0398$ under the observational design.  The misspecification is
identical in both cases --- the same wrong model, the same DGP --- so
the difference must come entirely from the assignment mechanism.

The DGP makes the cancellation under CRD explicit.  Under CRD,
$T \indep X$, so $X$ within each arm is still $\mathrm{Uniform}(0,1)$.
The population OLS line fit to $Y(t) = t + 3(1+t)X^2 + \varepsilon$
against $X$ over $U(0,1)$ has slope
$\beta_t = 3(1+t)\,{\mathrm{Cov}(X^2,X)}/{\mathrm{Var}(X)} = 3(1+t)$
(since both covariance and variance equal $\tfrac{1}{12}$ for
$X\sim U(0,1)$), giving
$\hat\mu_{\mathrm{lin}}(0,x) = -\tfrac{1}{2}+3x$ and
$\hat\mu_{\mathrm{lin}}(1,x) = 6x$.
The pointwise error in the estimated CATE is therefore
\[
  b(x)
  \;=\;
  \bigl[\hat\mu_{\mathrm{lin}}(1,x) - \hat\mu_{\mathrm{lin}}(0,x)\bigr]
  - \tau(x)
  \;=\;
  \Bigl(\tfrac{1}{2}+3x\Bigr) - \bigl(1+3x^2\bigr)
  \;=\;
  -\tfrac{1}{2}+3x-3x^2.
\]
The pointwise bias is nonzero everywhere: the linear model
misrepresents the curved surface at every $x$.  Yet its average over
$X\sim U(0,1)$ is exactly zero:
\[
  \E\bigl[b(X)\bigr]
  = \int_0^1\!\bigl(-\tfrac{1}{2}+3x-3x^2\bigr)\,dx
  = \Bigl[-\tfrac{x}{2}+\tfrac{3x^2}{2}-x^3\Bigr]_0^1
  = -\tfrac{1}{2}+\tfrac{3}{2}-1 = 0.
\]
The cancellation is structural, not accidental.  The population OLS
line is the linear projection of $\tau(X)$ onto $\{1, X\}$ under the
fitting distribution.  A linear projection always preserves the mean,
so $\E[\hat\tau_{\mathrm{lin}}(X)] = \E[\tau(X)] = \tau_{\mathrm{ATE}}$
whenever the expectation on the left is taken over the \emph{same}
distribution used for fitting.  Under CRD the fitting distribution
(the distribution of $X$ within each arm) equals the evaluation
distribution (the marginal of $X$), so that condition holds and the
OR estimator is consistent despite misspecification.

Under the observational design, the fitting and evaluation distributions
diverge.  Because $T = \mathrm{Bern}(\mathrm{expit}(-2+5X))$, high-$X$
units are almost always treated and low-$X$ units almost never.  The
model for arm $t=1$ is fit on a sample concentrated at high $X$, and
the model for arm $t=0$ on a sample concentrated at low $X$.  The
fitted lines are tuned to different parts of the covariate space, so
when both are evaluated at the full marginal of $X$, the pointwise
errors no longer cancel.  Model misspecification is harmless under
randomization; it becomes a source of bias precisely when confounding
is present.
\end{warn}

\begin{remark}[What This Lab Does Not Show]
\label{w5:rem:lab-limits}
The simulation fixes the bandwidth $h = 0.5 \cdot n^{-1/5}$ and does
not perform data-driven bandwidth selection.  In practice, choosing $h$
by cross-validation changes the bias--variance tradeoff and can
substantially improve the local-linear estimator's RMSE.  The lab also
studies only a scalar $X$; with a multivariate covariate, the curse of
dimensionality would make local linear regression impractical, which
is one motivation for the propensity score approach of Chapter~6.
\end{remark}

% ============================================================
\section{Where the Propensity Score Fits}
% ============================================================

The three methods developed in §§3--5 --- regression adjustment,
standardization, and stratification --- all implement the same
identifying formula~\eqref{w5:eq:backdoor}.  They differ in how they
deal with the conditioning on the high-dimensional covariate vector $X$:

\medskip
\begin{center}
\renewcommand{\arraystretch}{1.4}
\begin{tabular}{p{3.5cm} p{5cm} p{4.5cm}}
\toprule
\textbf{Method} & \textbf{How it handles $X$} & \textbf{Key assumption} \\
\midrule
Regression / standardization
  & Models $\E[Y \mid T,X]$ parametrically
  & Outcome model correctly specified \\[4pt]
Stratification
  & Groups units by $X$; estimates effect in each cell
  & Strata narrow enough for approximate balance \\[4pt]
Propensity score (Ch.~6)
  & Reduces $X$ to scalar $\pi(X) = P(T{=}1 \mid X)$
  & Propensity model correctly specified \\
\bottomrule
\end{tabular}
\end{center}
\renewcommand{\arraystretch}{1}

\medskip
When $X$ is low-dimensional, stratification and regression work well.
When $X$ has many components, direct stratification suffers from the
curse of dimensionality.  The \emph{propensity score} solves this by
proving (Theorem~\ref{w6:thm:ps}, Chapter~6) that adjusting for the
scalar $\pi(X)$ is \emph{sufficient} for identification, even when $X$ is
high-dimensional.

\begin{tcolorbox}[defstyle, title={Bridge to Chapter~6}]
Chapter~6 takes up this idea formally.  It proves the balancing property
($T \indep X \mid \pi(X)$), derives the Propensity Score Theorem
(strong ignorability given $X$ implies strong ignorability given $\pi(X)$
alone), and develops three estimation strategies based on the propensity
score: inverse probability weighting (IPW), matching, and propensity
score stratification.  Matching, in particular, need not rely on the
propensity score: covariate-distance methods such as Mahalanobis-distance
matching and nearest-neighbor matching on $X$ directly are equally valid,
and are introduced alongside propensity score matching in Chapter~6.
\end{tcolorbox}

% ============================================================
\section{Summary}
% ============================================================

\begin{enumerate}
  \item \textbf{Randomization as graph surgery.}  A randomized
    experiment removes all arrows into $T$, so the fundamental equality
    $f(y \mid \doop(T{=}t)) = f(y \mid T{=}t)$ holds and the
    difference-in-means estimator is unbiased for the ATE without
    covariate adjustment.

  \item \textbf{Ignorability is the key assumption.}  Under
    unconfoundedness $(Y(0),Y(1)) \indep T \mid X$, the back-door
    adjustment formula identifies the ATE.  This assumption is
    substantive, design-determined, and untestable from observational
    data.

  \item \textbf{Three estimation strategies.}  Regression adjustment
    (G-formula), standardization, and stratification all implement the
    same identifying formula via different modeling choices.  All
    require correct specification of either an outcome model (for
    regression/standardization) or a stratification scheme fine enough
    to achieve balance.

  \item \textbf{Simpson's paradox.}  Pooled associations can reverse
    within subgroups when a confounder determines treatment selection.
    The correct causal estimate requires standardization over the
    marginal covariate distribution, guided by the back-door criterion.

  \item \textbf{The propensity score dimension reduction.}  When $X$
    is high-dimensional, direct stratification or cell-by-cell
    standardization breaks down.  The propensity score $\pi(X)$ provides
    a scalar sufficient statistic for adjustment.  Chapter~6 develops
    the theory and estimation methods.
\end{enumerate}

% ============================================================
\section*{Problems}
\addcontentsline{toc}{section}{Problems}
% ============================================================

\begin{enumerate}

\item \textbf{Randomization and the do-calculus.}
  Let the observational DAG be
  $\{U \to T,\; U \to Y,\; X \to T,\; T \to Y\}$
  with $U$ unobserved.
  \begin{enumerate}
    \item Identify all back-door paths from $T$ to $Y$.
    \item Does $X$ satisfy the back-door criterion?  Does $U$?  Explain.
    \item Now suppose treatment is randomized.  Draw the modified DAG
      and identify all back-door paths.  Show that
      $f(y \mid \doop(T{=}t)) = f(y \mid T{=}t)$ holds in the
      randomized DAG using Rule~2 of the do-calculus.
    \item Under randomization, is covariate adjustment on $X$ necessary
      for unbiasedness?  Is it ever beneficial?  Explain.
  \end{enumerate}

\item \textbf{Ignorability and the back-door criterion.}
  Consider the DAG $\{X \to T,\; X \to Y,\; T \to Y,\; U \to Y\}$
  with $U$ unobserved.
  \begin{enumerate}
    \item Does $\{X\}$ satisfy the back-door criterion?  Write out the
      identifying formula for $\tau_{\mathrm{ATE}}$.
    \item Now add the edge $U \to T$ to the DAG.  Does $\{X\}$ still
      satisfy the back-door criterion?  What does the back-door criterion
      require when both $X$ and $U$ affect $T$?
    \item Explain in words what ``unconfoundedness given $X$'' means
      about the role of $U$ in the data-generating process.
  \end{enumerate}

\item \textbf{Standardization.}
  A study of job training ($T$) and earnings ($Y$, in thousands)
  produces the following cell means, with $P(X{=}0) = 0.4$ and
  $P(X{=}1) = 0.6$:
  \begin{center}
  \begin{tabular}{lccc}
  \toprule
  $X$ & $\hat\mu(1,x)$ & $\hat\mu(0,x)$ & $n_x/n$ \\
  \midrule
  0 & 28 & 22 & 0.4 \\
  1 & 35 & 31 & 0.6 \\
  \bottomrule
  \end{tabular}
  \end{center}
  \begin{enumerate}
    \item Compute $\hat\tau_{\mathrm{ATE}}$ using the standardization
      formula.
    \item Compute $\hat\tau_{\mathrm{ATT}}$, given that all treated
      units come from $X{=}1$.
    \item The unadjusted difference in means is $33 - 25 = 8$.  Compare
      to your answers in (a) and (b) and explain the discrepancy.
  \end{enumerate}

\item \textbf{Stratification and Cochran's rule.}
  You have a binary confounder $X \in \{0,1\}$ and form two strata.
  Within stratum $X{=}0$: $n_0 = 600$, $\bar Y_1 = 10$,
  $\bar Y_0 = 8$.  Within stratum $X{=}1$: $n_1 = 400$,
  $\bar Y_1 = 15$, $\bar Y_0 = 12$.
  \begin{enumerate}
    \item Compute the stratified ATE estimator $\hat\tau_{\mathrm{STRAT}}$.
    \item Suppose an unadjusted difference-in-means gives
      $\hat\tau_{\mathrm{DIM}} = 5.5$.  Explain why the two estimates
      differ and which one is the appropriate causal estimate.
    \item Describe one limitation that would arise if $X$ were a
      continuous variable with 15 dimensions.
  \end{enumerate}

\item \textbf{Simpson's paradox.}
  A hospital reports that patients who receive intensive care (ICU
  admission, $T{=}1$) have a higher mortality rate than those who do
  not ($T{=}0$): $P(Y{=}1 \mid T{=}1) = 0.30$ vs.\
  $P(Y{=}1 \mid T{=}0) = 0.10$.
  \begin{enumerate}
    \item Construct a numerical example (a $2 \times 2 \times 2$
      table with disease severity $X \in \{\text{mild}, \text{severe}\}$
      as the confounder) that is consistent with the above pooled
      numbers and yet shows ICU admission \emph{reduces} mortality
      within both severity strata.
    \item Compute the standardized ATE using the marginal distribution
      of $X$.
    \item Draw the DAG for this example.  Identify the back-door path
      that the pooled comparison fails to block.
    \item A colleague suggests that since the hospital treats patients
      with both mild and severe illness, the pooled statistic is the
      right answer.  Explain, using potential outcomes notation, why
      this argument is incorrect.
  \end{enumerate}

\end{enumerate}

% Bibliography (standalone only)
\ifSubfilesClassLoaded{%
  \bibliographystyle{plainnat}
  \bibliography{causal_inference}
}{}

\end{document}